\chapter{Complexity — Universal}
%%% generated by the blueprint pipeline from TCSlib.Complexity.CircuitComplexity.Universal
\section{Overview}

Declarations of the Lean module \texttt{TCSlib.Complexity.CircuitComplexity.Universal}.

\section{Declarations}

\begin{theorem}[Evaluation of a Boolean literal]
\lean{BoolCircuit.lit_eval_iff}
\label{BoolCircuit.lit_eval_iff}
\leanok
\uses{BoolCircuit.Lit.eval}
\difficulty{4}
Let $n$ be a natural number, let $i$ be an index among $n$ variables, let
$b\in\{\mathrm{false},\mathrm{true}\}$, and let $x$ be a Boolean assignment to those
variables. The literal at index $i$ with polarity $b$ evaluates to true if and only if
$x(i)=b$.
\end{theorem}

\begin{theorem}[Size of a list of equal-size circuits]
\lean{BoolCircuit.foldr_size_map}
\label{BoolCircuit.foldr_size_map}
\leanok
\uses{BoolCircuit.TreeCircuit, BoolCircuit.TreeCircuit.size}
\difficulty{4}
Let $n$ and $k$ be natural numbers, let $g$ map a type $\alpha$ to Boolean circuit trees
on $n$ variables, and suppose every circuit $g(a)$ has size $k$. Then the sum of the
sizes of the circuits $g(a)$ over any finite list $l$ of elements of $\alpha$ is $|l|k$.
\end{theorem}

\begin{definition}[Minterm circuit]
\lean{BoolCircuit.minterm}
\label{BoolCircuit.minterm}
\leanok
\uses{BoolCircuit.TreeCircuit}
For a natural number $n$ and a Boolean assignment $v$ on the $n$ variables, form an AND
gate whose children, in index order from $0$ to $n-1$, are the literals at those indices
with polarity given by $v$.
\end{definition}

\begin{theorem}[Minterm characterizes an assignment]
\lean{BoolCircuit.minterm_eval_iff}
\label{BoolCircuit.minterm_eval_iff}
\leanok
\uses{BoolCircuit.TreeCircuit.eval, BoolCircuit.TreeCircuit.eval_lit, BoolCircuit.TreeCircuit.eval_node_true_iff, BoolCircuit.lit_eval_iff, BoolCircuit.minterm}
\difficulty{5}
Let $n$ be a nonnegative integer, and let $v$ and $x$ be Boolean assignments on $n$
variables. The minterm specified by $v$ evaluates to true on $x$ if and only if $x=v$.
\end{theorem}

\begin{theorem}[Size of a minterm]
\lean{BoolCircuit.minterm_size}
\label{BoolCircuit.minterm_size}
\leanok
\uses{BoolCircuit.TreeCircuit.size, BoolCircuit.foldr_size_map, BoolCircuit.minterm}
\difficulty{3}
For any nonnegative integer $n$ and Boolean assignment $v$ on $n$ variables, the minterm
specified by $v$ has size $n+1$.
\end{theorem}

\begin{definition}[Universal Boolean circuit]
\lean{BoolCircuit.universalCircuit}
\label{BoolCircuit.universalCircuit}
\leanok
\uses{BoolCircuit.TreeCircuit, BoolCircuit.minterm}
For a Boolean function $f$ on $n$ Boolean variables, form a tree circuit whose root is
an OR gate and whose children are the minterms corresponding to exactly those inputs $v$
for which $f(v)=1$.
\end{definition}

\begin{theorem}[Evaluation of the Universal Circuit]
\lean{BoolCircuit.universalCircuit_eval}
\label{BoolCircuit.universalCircuit_eval}
\leanok
\uses{BoolCircuit.TreeCircuit.eval, BoolCircuit.TreeCircuit.eval_node_false_iff, BoolCircuit.minterm, BoolCircuit.minterm_eval_iff, BoolCircuit.universalCircuit}
\difficulty{4}
Let $n$ be a natural number, let $f$ be a Boolean function on assignments to $n$
variables, and let $x$ be such an assignment. The universal circuit representing $f$,
evaluated at $x$, has value $f(x)$.
\end{theorem}

\begin{theorem}[Size of the Universal Circuit]
\lean{BoolCircuit.universalCircuit_size}
\label{BoolCircuit.universalCircuit_size}
\leanok
\uses{BoolCircuit.TreeCircuit.size, BoolCircuit.foldr_size_map, BoolCircuit.minterm, BoolCircuit.minterm_size, BoolCircuit.universalCircuit}
\difficulty{3}
Let $n$ be a natural number and $f$ a Boolean function on assignments to $n$ variables.
The universal circuit representing $f$ has size $(n+1)$ times the number of assignments
on which $f$ is true, plus one.
\end{theorem}

\begin{theorem}[Upper Bound on Universal Circuit Size]
\lean{BoolCircuit.universalCircuit_size_le}
\label{BoolCircuit.universalCircuit_size_le}
\leanok
\uses{BoolCircuit.TreeCircuit.size, BoolCircuit.universalCircuit, BoolCircuit.universalCircuit_size}
\difficulty{3}
Let $n$ be a natural number and $f$ a Boolean function on assignments to $n$ variables.
The universal circuit representing $f$ has size at most $2^n(n+1)+1$.
\end{theorem}

\begin{theorem}[Small tree circuit for every Boolean function]
\lean{BoolCircuit.exists_circuit_eval_eq_size_le}
\label{BoolCircuit.exists_circuit_eval_eq_size_le}
\leanok
\uses{BoolCircuit.TreeCircuit, BoolCircuit.TreeCircuit.eval, BoolCircuit.TreeCircuit.size, BoolCircuit.universalCircuit, BoolCircuit.universalCircuit_eval, BoolCircuit.universalCircuit_size_le}
\difficulty{1}
For every nonnegative integer $n$ and every Boolean function on assignments to $n$
variables, there is a Boolean circuit tree that computes the function on every
assignment and has size at most $2^n(n+1)+1$.
\end{theorem}

\begin{theorem}[Universal Circuit for the Zero Function]
\lean{BoolCircuit.universalCircuit_const_false}
\label{BoolCircuit.universalCircuit_const_false}
\leanok
\uses{BoolCircuit.universalCircuit}
\difficulty{1}
For every natural number $n$, the universal circuit representing the constant-zero
Boolean function on $n$ variables is an OR gate with no children.
\end{theorem}

\begin{theorem}[Size of the Universal Circuit for the One Function]
\lean{BoolCircuit.universalCircuit_const_true_size}
\label{BoolCircuit.universalCircuit_const_true_size}
\leanok
\uses{BoolCircuit.TreeCircuit.size, BoolCircuit.universalCircuit, BoolCircuit.universalCircuit_size}
\difficulty{2}
For every natural number $n$, the universal circuit representing the constant-one
Boolean function on $n$ variables has size $2^n(n+1)+1$.
\end{theorem}

\begin{theorem}[Evaluation of the zero-variable minterm]
\lean{BoolCircuit.minterm_eval_zero}
\label{BoolCircuit.minterm_eval_zero}
\leanok
\uses{BoolCircuit.TreeCircuit.eval, BoolCircuit.minterm, BoolCircuit.minterm_eval_iff}
\difficulty{2}
For the unique Boolean assignments $v$ and $x$ on zero variables, the minterm specified
by $v$ evaluates to true on $x$.
\end{theorem}
